% !LOOM digest: manolache_VirtualPullbacks2012
% !LOOM prefix: manolacheVirtualPullbacks2012
% !LOOM extracted-from: doi:10.1090/S1056-3911-2011-00606-1
% !LOOM method: extract
% !LOOM created: 2026-09-18
% !LOOM requires: latexsym, mathrsfs, xy

\section*{Overview}
Given a regular embedding of schemes $i:X\to Y$ one can construct a ``well-behaved'' morphism  (i.e. a bivariant class) $i^!:A_*(Y)\to A_*(X)$ by applying the Fulton-MacPherson construction (see \cite{f}, Chapter 6). The construction of this morphism can be generalized to the following setting. If $i:X\to Y$ is a closed embedding of schemes and $C_{X/Y}\subset E$ is a closed embedding of the normal cone of $i$ into a rank-$r$ vector bundle, then these data determine a ``well-behaved'' morphism $i^!:A_*(Y)\to A_{(*-r)}(X)$ (see Example 17.6.4 in \cite{f}).
\\The main result of this paper generalizes Fulton's example to stacks and to a larger class of morphisms. These generalizations allow us to view virtual classes (see \cite{lt}, \cite{bf}) as generalized pull-backs. This approach will allow us to deduce certain relations between Gromov-Witten invariants, which were the initial reasons for doing this work.
\\It should be said that the idea is not entirely new, although we did not find this approach in the literature. The main inspiration point was the ``functoriality property of the Behrend-Fantechi class'' of Kim, Kresch and Pantev in \cite{K:01}. Also, a similar situation appears in Jun Li's $[M,N]^{\mathrm{virt}}$-construction (see \cite{Li:01}).
\\When this paper was in an advanced state, we have been informed of Hsin-Hong Lai's paper on ``Gromov-Witten invariants of blow-ups along manifolds with convex normal bundle'' (\cite{h}). The ideas there are basically the same, with a slightly different flavor. We hope, however, that our point of view will contribute to a clear understanding of this subject.
\\
\\In the first section we recall the notions of normal cones of Behrend-Fantechi and Kresch and prove that these two notions are canonically isomorphic. This allows us to use Kresch's ``deformation to the normal cone'' in the context of ``intrinsic normal cones''.
\\The main idea of the second section is to replace the normal sheaf $N_{F/G}$ of a morphism of stacks $f:F\to G$ with a ``virtual normal bundle''. The appropriate context for this is given by obstruction theories. Precisely, if $f$ is a DM-type morphism of Artin stacks (see Definition \ref{manolacheVirtualPullbacks2012-def-2.1}) which admits a perfect relative obstruction theory $E_{F/G}^{\bullet}$ (see \cite{bf}), then we take the virtual normal bundle to be $h^1/h^0((E_{F/G}^{\vee})^{\bullet})$. Using this, we obtain a well-defined morphism $f^!:A_*(G)\to A_*(F)$, that we call a virtual pull-back. As a byproduct of our construction we obtain a generalized notion of virtual fundamental class which applies to some examples of Artin stacks.
\\In Section 3 we show that the virtual pull-back satisfies the usual compatibility conditions. Moreover, when we deal with stacks possessing virtual classes we prove that subject to a very natural compatibility relation between obstructions (see Definition \ref{manolacheVirtualPullbacks2012-def-4.5}) the construction gives a map that sends the virtual class of $G$ to the virtual class of $F$. The statement may also be seen as a generalization of the functoriality property in \cite{bf} and \cite{K:01}.
\\As an application we provide the answer to a very natural question. Given a smooth projective variety $X$ and its blow-up $p:\tilde{X}\to X$ along some smooth projective subvariety, we would like to know when do certain Gromov-Witten invariants of $X$ and $\tilde{X}$ agree. More precisely, if we start with a given homology class $\beta\in A_1(X)$ and a collection of cohomology classes $\gamma_i\in A^*(X)$, then we can associate a ``lifted'' homology class in $A_1(\tilde{X})$ (see Definition \ref{manolacheVirtualPullbacks2012-def-5.10} for a precise statement) and cohomology classes $p^*\gamma_i\in A^*(\tilde{X})$. One could expect that the Gromov-Witten invariants associated to these data should be equal. This was first analyzed by Gathmann (\cite{G:01}) where $X$ was a convex space and $Y$ a point and by Hu (\cite{H:00}, \cite{H:02}) where it was treated the blow-up along points, curves and surfaces. Recently, it was shown and by Lai (\cite{h}) that (subject to a minor condition) the expectation is true  for genus zero Gromov-Witten invariants of blow-ups along subvarieties with convex normal bundles. Our idea is to show the equality of rational Gromov-Witten invariants for $X$ convex and then ``pull the relation back'' to an arbitrary $X$ (see Proposition \ref{manolacheVirtualPullbacks2012-prop-5.14}). The statement we get should be compared with Theorem 1.6 in \cite{h}.
\\The second application concerns rational Gromov-Witten invariants of projective bundles $p_X:\mathbb{P}_X(V)\to X$. These were studied by Ruan and Qin (\cite{ru}) where $X$ was taken to be a projective space and by Elezi (\cite{e1}, \cite{e2}) when $V$ is a split bundle and $X$ is a toric variety. Here, we analyze the map induced by $p_X$ between the corresponding moduli spaces of stable maps to $\mathbb{P}_X(V)$ and $X$.
\\We also show that a particular case of Costello's push-forward formula follows as an easy consequence of our formalism.
\paragraph{Notation and conventions.}
We work over a fixed ground field.
\\An Artin stack is an algebraic stack in the sense of \cite{lau} of finite type over the ground field. For simplicity we will call Deligne --Mumford stacks DM stacks.
\\Unless otherwise specified we will try to respect the following convention: we will usually denote schemes by $X,\ Y,\ Z$, etc, Artin stacks by $F,\ G,\ H$, etc. and Artin stacks for which we know that they are not Deligne-Mumford stacks (such as the moduli space of genus-$g$ curves or vector bundle stacks) by gothic letters $\mathfrak{M}_{g},\ \mathfrak{E},\ \mathfrak{F}$, etc.
\\By a commutative diagram of stacks we mean a 2-commutative diagram of stacks and by a cartesian diagram of stacks we mean a 2-cartesian diagram of stacks.
\\Chow groups for schemes are defined in the sense of \cite{f}; this definition has been extended to DM stacks (with $\mathbb{Q}$-coefficients) by Vistoli (\cite{v}) and to algebraic stacks (with $\mathbb{Z}$-coefficients) by Kresch (\cite{k}). We will consider Chow groups (of schemes/stacks) with $\mathbb{Q}$-coefficients.
\\For a fixed stack $F$ we denote by $\mathcal{D}_F$ the derived category of quasicoherent $\mathcal{O}_F$ modules. Unless otherwise specified, we will denote the derived functors ${\bf{L}}f^*$, ${\bf{R}}f_*$, ${\bf{R}}Hom$, etc. by $f^*$, $f_*$, $Hom$, etc.
\\For a fixed stack $F$ we denote by $L_F$ its cotangent complex defined in \cite{ol}.

\paragraph{Acknowledgments.} I would like to thank Dan Abramovich, Claudio Fontanari, Lothar G\"{o}ttsche, Andrew Kresch, Hsin-Hong Lai and Ravi Vakil for useful discussions.
\\I am grateful to my advisors Ionu\c{t} Ciocan-Fontanine and Barbara Fantechi for helpful discussions and suggestions. Special thanks are due to Ionu\c{t} for many helpful conversations and hospitality during a one month stay at the University of Minnesota. I am mostly grateful to Barbara and the anonymous referee for their huge amount of suggestions which led to a major improvement of the exposure.
\\I learned a lot during the Moduli Spaces program at the Mittag-Leffler Institute and at the University of Minnesota. It is a pleasure to thank both institutions for the wonderful research environment and support.
\\I was partially supported by an ENIGMA grant.

\begin{theorem}[{\cite[Standing assumptions]{manolache_VirtualPullbacks2012}}]\label{manolacheVirtualPullbacks2012-setup}
\incomplete{Standing assumptions not extracted; see the paper.}
\end{theorem}

\section{Preliminaries}\label{manolacheVirtualPullbacks2012-sec-2}

\subsection{Intrinsic normal cones to DM-type morphisms}\label{manolacheVirtualPullbacks2012-sec-2.1}

\begin{definition}[{\cite[Definition 2.1, p.~4]{manolache_VirtualPullbacks2012}}]\label{manolacheVirtualPullbacks2012-def-2.1}
\label{manolacheVirtualPullbacks2012-dm}
 A morphism $f:F\to G$ of Artin stacks is called of Deligne-Mumford type (or shortly of DM-type) if for any morphism $V\to G$, with $V$ a scheme, $F\times_GV$ is a Deligne-Mumford stack.
\end{definition}

\begin{remark}[{\cite[Remark 2.2]{manolache_VirtualPullbacks2012}}]\label{manolacheVirtualPullbacks2012-rem-2.2}
Let us consider the following Cartesian diagram
\begin{equation}
 \xymatrix {F^{\prime} \ar[r]^{f^{\prime}}\ar[d] & G^{\prime}\ar[d]\\
             F\ar[r]^f& G.}
\end{equation}
If $f$ is a DM-type morphism, then $f^{\prime}$ is a DM-type morphism.
\end{remark}

\begin{remark}[{\cite[Remark 2.3]{manolache_VirtualPullbacks2012}}]\label{manolacheVirtualPullbacks2012-rem-2.3}
Let $f:F\to G$ be morphism of stacks and let $L_f$ be the relative cotangent complex. Then $f$ is of DM-type if and only if $L_f\in \mathcal{D}^{\leq0}_F$
\end{remark}

\begin{definition}[{\cite[Definition 2.4]{manolache_VirtualPullbacks2012}}]\label{manolacheVirtualPullbacks2012-def-2.4}
Let $E^{\bullet}\in\mathcal{D}^{\leq0}_F$. $E^{\bullet}$ is said to be of perfect amplitude if there exists $n\geq0$ such that $E^{\bullet}$ is locally isomorphic to $[E^{-n}\to...\to E^{0}]$, where $\forall i\in\{-n,...,0\}$, $E^i$ is a locally free sheaf.
\end{definition}

\begin{definition}[{\cite[Definition 2.5]{manolache_VirtualPullbacks2012}}]\label{manolacheVirtualPullbacks2012-def-2.5}
Let $E^{\bullet}\in\mathcal{D}^{\leq0}_X$. Then a homomorphism $\Phi: E^{\bullet}\to L^{\bullet}$ in $\mathcal{D}_F$ is called an obstruction theory if $h^0(\Phi)$ is an isomorphism and $h^{-1}(\Phi)$ is surjective. If moreover, $E^{\bullet}$ is of perfect amplitude, then $E^{\bullet}$ is called a perfect obstruction theory.
\end{definition}

\begin{convention}[{\cite[Convention 2.6]{manolache_VirtualPullbacks2012}}]\label{manolacheVirtualPullbacks2012-convention-2.6}
Unless otherwise stated by a perfect obstruction theory we will always mean of perfect amplitude  contained in $[-1,0]$.
\end{convention}

\begin{definition}[{\cite[Definition 2.7]{manolache_VirtualPullbacks2012}}]\label{manolacheVirtualPullbacks2012-def-2.7}
Let $X$ be a scheme and $\mathcal{F}$ be a coherent sheaf on $X$. We call $C(\mathcal{F}):=Spec Sym(\mathcal{F})$ an abelian cone over $X$.
\end{definition}

\begin{definition}[{\cite[Definition 2.8]{manolache_VirtualPullbacks2012}}]\label{manolacheVirtualPullbacks2012-def-2.8}
 An $\mathbb{A}^1$-invariant subscheme of $C(\mathcal{F})$ that contains the zero section is called a cone over $X$.
\end{definition}

\begin{definition}[{\cite[Definition 2.9]{manolache_VirtualPullbacks2012}}]\label{manolacheVirtualPullbacks2012-def-2.9}
Let $F$ be a stack and let $E^{\bullet}:=[E^0\to E^1]$ be an element in $\mathcal{D}_F$ such that $E^{i}$ is an abelian sheaf for $i=0,1$ and $h^i({E^{\bullet}})$ is coherent for $i=0,1$. We call the stack quotient $[E^1/E^0]$ (in the sense of \cite{bf} Section 2) an abelian cone stack over stack $F$ .
\\A cone stack is a closed substack of an abelian cone stack invariant under the action of $\mathbb{A}^1$ and containing the zero section.
\end{definition}

\begin{convention}[{\cite[Convention 2.10, p.~6]{manolache_VirtualPullbacks2012}}]\label{manolacheVirtualPullbacks2012-convention-2.10}
\label{manolacheVirtualPullbacks2012-stack}From now on, unless otherwise stated, by cones we will mean cone-stacks.
\end{convention}

\begin{example}[{\cite[Example 2.11]{manolache_VirtualPullbacks2012}}]\label{manolacheVirtualPullbacks2012-ex-2.11}
(i) Let $i:X\to Y$ be a closed embedding of schemes. If $\mathcal{I}$ denotes the ideal sheaf of $X$ in $Y$, then $N_{X/Y}=Spec Sym\ \mathcal{I}/\mathcal{I}^2$ is called the normal sheaf of $X$ in $Y$ and $C_{X/Y}:=Spec\oplus_{k\geq 0} \mathcal{I}^k/\mathcal{I}^{k+1}\hookrightarrow N_{X/Y}$ is called the normal cone of $X$ in $Y$.
\\(ii) If $f:F\to G$ is a local immersion of DM-stacks, then Vistoli defines (see \cite{v}, Definition 1.20) the normal cone to $f$ as described below. Let us consider a commutative diagram
\begin{equation}\label{manolacheVirtualPullbacks2012-vi}
 \xymatrix {U \ar[r]^{\tilde{f}}\ar[d] & V\ar[d]\\
             F\ar[r]^f& G}
\end{equation}
with $U$, $V$ schemes, the upper horizontal arrow a closed immersion and the vertical arrows \'{e}tale. Then $C_{F/G}$ is the cone obtained by descent from $C_{U/V}$.
\\Note that $C_{F/G}\hookrightarrow N_{F/G}=SpecSym\ h^{-1}(L_{F/G})$.
\end{example}

\begin{definition}[{\cite[Definition 2.12, p.~6]{manolache_VirtualPullbacks2012}}]\label{manolacheVirtualPullbacks2012-def-2.12}
\label{manolacheVirtualPullbacks2012-normsheaf}
 Let $f:F\to G$ be a DM-type morphism and let $L_{F/G}\in$ ob $\mathcal{D}(\mathcal{O}_{F})$ be the cotangent complex. Then we denote the stack $$h^1/h^0(L_{F/G}^{\vee}):=h^1/h^0(\tau_{[0,1]}(L_{F/G}^{\vee}))$$ (see \cite{bf}) by $\mathfrak{N}_{F/G}$ and we call it the \emph{intrinsic normal sheaf}.
\end{definition}

\begin{proposition}[{\cite[Proposition 2.13]{manolache_VirtualPullbacks2012}}]\label{manolacheVirtualPullbacks2012-prop-2.13}
(Behrend-Fantechi) Let us consider diagram (\ref{manolacheVirtualPullbacks2012-vi}) with $F$ a Deligne-Mumford stack, the upper horizontal arrow a closed immersion, $U\to F$ an \'{e}tale morphism and $V\to G$ a smooth morphism. Then for any $U$ and $V$ as above, there exists a unique cone-stack $\mathfrak{C}_{F/G}\subset\mathfrak{N}_{F/G}$ such that $\mathfrak{C}_{F/G}\times_FU=[C_{U/V}/\tilde{f}^*T_{V/G}]$.
\end{proposition}

\begin{definition}[{\cite[Definition 2.14, p.~6]{manolache_VirtualPullbacks2012}}]\label{manolacheVirtualPullbacks2012-def-2.14}
\label{manolacheVirtualPullbacks2012-befa}We call $\mathfrak{C}_{F/G}$ the \emph{intrinsic normal cone to $f$}.
\end{definition}

\begin{remark}[{\cite[Remark 2.15]{manolache_VirtualPullbacks2012}}]\label{manolacheVirtualPullbacks2012-rem-2.15}
 The notion of intrinsic normal cone to a morphism $F\to G$ from a Deligne-Mumford stack to an Artin stack has been introduced in \cite{bf} under the hypothesis $G$ is smooth. However, this hypothesis is not used in the construction of the cone. This restriction is only needed in \cite{bf} in order to define a virtual class (see \ref{manolacheVirtualPullbacks2012-cor-3.12} in the next section).
\end{remark}

\begin{lemma}[{\cite[Lemma 2.16, p.~7]{manolache_VirtualPullbacks2012}}]\label{manolacheVirtualPullbacks2012-lem-2.16}
\label{manolacheVirtualPullbacks2012-seqcones}
 Let \begin{equation*}
 \xymatrix {F^{\prime} \ar[r]^{f^{\prime}}\ar[d]_{p} & G^{\prime}\ar[d]^q\\
             F\ar[r]^f & G}
\end{equation*}
be a commutative diagram of Artin stacks with $f$ and $f^{\prime}$ of DM-type. Then, the natural morphism $p^*L_{F/G}\to L_{F^{\prime}/G^{\prime}}$ in \cite{ol} induces morphism of abelian cone stacks $$\beta:\mathfrak{N}_{F^{\prime}/G^{\prime}}\rightarrow p^*\mathfrak{N}_{F/G}.$$
\end{lemma}

\begin{remark}[{\cite[Remark 2.17, p.~7]{manolache_VirtualPullbacks2012}}]\label{manolacheVirtualPullbacks2012-rem-2.17}
\label{manolacheVirtualPullbacks2012-etale} In notations as in \ref{manolacheVirtualPullbacks2012-lem-2.16} let $p:F^{\prime}\to F$ be an \'etale morphism, $G^{\prime}=G$ and $q$ the identity morphism. Then, $\beta:\mathfrak{N}_{F^{\prime}/G}\rightarrow p^*\mathfrak{N}_{F/G}$ is an isomorphism. To see this let us consider the distinguished triangle $$p^*L_{F/G}\to L_{F^{\prime}/F}\to L_{F^{\prime}/G}\to p^*L_{F/G}[1].$$ The claim follows from the fact that $L_{F^{\prime}/F}=0$.
\end{remark}

\begin{proposition}[{\cite[Proposition 2.18, p.~7]{manolache_VirtualPullbacks2012}}]\label{manolacheVirtualPullbacks2012-prop-2.18}
\uses{manolacheVirtualPullbacks2012-lem-2.16}
\label{manolacheVirtualPullbacks2012-dmincl}
 Let \begin{equation*}
 \xymatrix {F^{\prime} \ar[r]^{f^{\prime}}\ar[d]_{p} & G^{\prime}\ar[d]^q\\
             F\ar[r]^f & G}
\end{equation*}
be a commutative diagram with $F$, $F^{\prime}$ DM-stacks stacks. Then, the morphism of Lemma \ref{manolacheVirtualPullbacks2012-lem-2.16} induces a morphism of cone stacks $\alpha:\mathfrak{C}_{F^{\prime}/G^{\prime}}\rightarrow p^*\mathfrak{C}_{F/G}$. If the diagram is cartesian, then $\alpha$ is a closed immersion. If moreover, $q$ is flat, then $\alpha$ is an isomorphism.
\end{proposition}

\begin{remark}[{\cite[Remark 2.19, p.~9]{manolache_VirtualPullbacks2012}}]\label{manolacheVirtualPullbacks2012-rem-2.19}
\label{manolacheVirtualPullbacks2012-injsurj}
Let us consider diagram 5. By  \cite{bf} a morphism of cones $C_{U^{\prime}/V^{\prime}}\to s^*C_{U/V}$ is a closed immersion if and only if the induced morphism $N_{U^{\prime}/V^{\prime}}\to s^*N_{U/V}$ is a closed immersion. This shows that $\mathfrak{C}_{U^{\prime}/G^{\prime}}\rightarrow p^*\mathfrak{C}_{U/G}$ is a closed immersion if and only if $\mathfrak{N}_{U^{\prime}/G^{\prime}}\rightarrow p^*\mathfrak{N}_{U/G}$ is a closed immersion. As being a closed immersion can be checked (\'etale) locally we see that $\mathfrak{C}_{F^{\prime}/G^{\prime}}\rightarrow p^*\mathfrak{C}_{F/G}$ is a closed immersion if and only if $\mathfrak{N}_{F^{\prime}/G^{\prime}}\rightarrow p^*\mathfrak{N}_{F/G}$ is a closed immersion.
\end{remark}

\begin{construct}[{\cite[Construction 2.20, p.~9]{manolache_VirtualPullbacks2012}}]\label{manolacheVirtualPullbacks2012-constr-2.20}
\label{manolacheVirtualPullbacks2012-locart} Let us consider the following commutative diagram
\begin{equation}\label{manolacheVirtualPullbacks2012-artincones}
 \xymatrix {U\ar[d]_r\ar[rd]\\
             F\ar[r]^f& G}
\end{equation}
with $f$ a DM-type morphism, $U$ a scheme and $r$ a smooth morphism. By Lemma \ref{manolacheVirtualPullbacks2012-lem-2.16} we have a morphism
\begin{equation}
\gamma:\mathfrak{N}_{U/G}\to r^*\mathfrak{N}_{F/G}.
\end{equation}
Let us consider the restriction $\tilde{\gamma}:\mathfrak{C}_{U/G}\to r^*\mathfrak{N}_{F/G}$. We denote the image of $\tilde{\gamma}$ by $\mathfrak{C}_{F/G}|_U$ and we call it the local normal cone on $U$ of $F$ to $G$ .
\end{construct}

\begin{lemma}[{\cite[Lemma 2.21]{manolache_VirtualPullbacks2012}}]\label{manolacheVirtualPullbacks2012-lem-2.21}
 Let us consider the following commutative diagram
\begin{equation*}
 \xymatrix {U^{\prime}\ar[d]_{r^{\prime}}\ar[rdd]\\U\ar[d]_r\ar[rd]\\
             F\ar[r]^f& G}
\end{equation*}
with $r$ and $r^{\prime}$ smooth morphisms. Then $\mathfrak{C}_{F/G}|_U\times_UU^{\prime}$ is naturally isomorphic to $\mathfrak{C}_{F/G}|_{U^{\prime}}$.
\end{lemma}

\begin{definition}[{\cite[Definition 2.22, p.~10]{manolache_VirtualPullbacks2012}}]\label{manolacheVirtualPullbacks2012-def-2.22}
\label{manolacheVirtualPullbacks2012-coneartinstack}
The cone $C_{F/G}$ is called the intrinsic normal cone of $f$, or when there is no risk of confusion the intrinsic normal cone of $F$ to $G$.
\end{definition}

\begin{remark}[{\cite[Remark 2.23]{manolache_VirtualPullbacks2012}}]\label{manolacheVirtualPullbacks2012-rem-2.23}
Let us consider diagram (\ref{manolacheVirtualPullbacks2012-artincones}) with $F$ a DM-stack and $r:U\to F$ an \'etale morphism.  By restricting the natural isomorphism $\mathfrak{N}_{U/G}\to r^*\mathfrak{N}_{F/G}$ to $\mathfrak{C}_{U/G}$ (see Remark \ref{manolacheVirtualPullbacks2012-rem-2.17}) we obtain an isomorphism of cones $\mathfrak{C}_{U/G}\to r^*\mathfrak{C}_{F/G}$. This shows that when $F$ is a DM-stack $C_f$ of Definition \ref{manolacheVirtualPullbacks2012-def-2.22} is canonically isomorphic to $C_f$ of Definition \ref{manolacheVirtualPullbacks2012-def-2.14}.
\end{remark}

\begin{remark}[{\cite[Remark 2.24]{manolache_VirtualPullbacks2012}}]\label{manolacheVirtualPullbacks2012-rem-2.24}
 Let us consider diagram (\ref{manolacheVirtualPullbacks2012-artincones}), with $r$ a smooth morphism. By Proposition \ref{manolacheVirtualPullbacks2012-prop-2.18} and Remark \ref{manolacheVirtualPullbacks2012-rem-2.19} we obtain a closed embedding $\mathfrak{C}_{U/F}\to\mathfrak{C}_{U/G}$. This shows that the sequence $0\to\mathfrak{N}_{U/F}\to\mathfrak{N}_{U/G}\to\mathfrak{N}_{F/G}|_U\to 0$ induces the exact sequence of cones (in the sense of Definition 1.12 in \cite{bf})
\begin{equation}\label{manolacheVirtualPullbacks2012-artco}
0\to\mathfrak{N}_{U/F}\to\mathfrak{C}_{U/G}\to\mathfrak{C}_{F/G}|_U\to 0.
\end{equation}
\end{remark}

\begin{remark}[{\cite[Remark 2.25]{manolache_VirtualPullbacks2012}}]\label{manolacheVirtualPullbacks2012-rem-2.25}
 Alternatively to Construction \ref{manolacheVirtualPullbacks2012-constr-2.20} one can define $\mathfrak{C}_{F/G}|_U$ by the sequence \ref{manolacheVirtualPullbacks2012-artco}.
\end{remark}

\begin{proposition}[{\cite[Proposition 2.26, p.~11]{manolache_VirtualPullbacks2012}}]\label{manolacheVirtualPullbacks2012-prop-2.26}
\uses{manolacheVirtualPullbacks2012-prop-2.18}
\label{manolacheVirtualPullbacks2012-pullbackofcones}
 Let \begin{equation*}
 \xymatrix {F^{\prime} \ar[r]^{f^{\prime}}\ar[d]_{p} & G^{\prime}\ar[d]^q\\
             F\ar[r]^f & G}
\end{equation*}
be a commutative diagram of Artin stacks with $f$ of DM-type. Then, there is an induced morphism of cone stacks $\alpha:\mathfrak{C}_{F^{\prime}/G^{\prime}}\rightarrow p^*\mathfrak{C}_{F/G}$. If moreover, the diagram is cartesian, then $\alpha$ is a closed immersion. If $q$ is flat, then $\alpha$ is an isomorphism.
\end{proposition}

\subsection{Normal cones to DM-type morphisms}\label{manolacheVirtualPullbacks2012-sec-2.2}

\begin{lemma}[{\cite[Lemma 2.27, p.~12]{manolache_VirtualPullbacks2012}}]\label{manolacheVirtualPullbacks2012-lem-2.27}
\label{manolacheVirtualPullbacks2012-reducetoincl}Let $f:F\to G$ be a DM-type morphism of Artin stacks. Then one can construct a commutative diagram (not unique)
\begin{equation}
 \xymatrix {U \ar[r]^{\tilde{f}}\ar[d]& V\ar[d]\\
             F\ar[r]^f& G}
\end{equation}
where $U$ and $V$ are schemes, the vertical arrows are smooth surjective and the top arrow $U\to V$ is  a closed immersion.
\end{lemma}

\begin{lemma}[{\cite[Lemma 2.28]{manolache_VirtualPullbacks2012}}]\label{manolacheVirtualPullbacks2012-lem-2.28}
\uses{manolacheVirtualPullbacks2012-prop-2.26}
Let $R:=U\times_{F}U$ and $S:=V\times_{G}V$, where $U$ and $V$ are defined in the proof of the previous Lemma. Then the natural map $R\to S$ is a locally closed immersion.
\end{lemma}

\begin{proposition}[{\cite[Proposition 2.29]{manolache_VirtualPullbacks2012}}]\label{manolacheVirtualPullbacks2012-prop-2.29}
 (Kresch) Let us consider the cone $C_{R/S}$. There are natural morphisms making $C_{R/S}\rightrightarrows C_{U/V}$ into a smooth groupoid in the category of schemes.
\end{proposition}

\begin{definition}[{\cite[Definition 2.30, p.~13]{manolache_VirtualPullbacks2012}}]\label{manolacheVirtualPullbacks2012-def-2.30}
\label{manolacheVirtualPullbacks2012-kresch}
 Let $C_{F/G}$ be the stack associated to the groupoid $[C_{R/S}\rightrightarrows C_{U/V}]$ and $N_{F/G}$ the stack associated to the groupoid $[N_{R/S}\rightrightarrows N_{U/V}]$. We call $C_{F/G}$ the normal cone of $f$ and $N_{F/G}$ the normal sheaf of $f$.
\end{definition}

\begin{theorem}[{\cite[Theorem 2.31, p.~13]{manolache_VirtualPullbacks2012}}]\label{manolacheVirtualPullbacks2012-thm-2.31}
\uses{manolacheVirtualPullbacks2012-lem-2.27}
\label{manolacheVirtualPullbacks2012-family} (Kresch)
 Let $f:F\to G$ be a DM-type morphism of Artin stacks. One can define a deformation space, i.e. a flat morphism $M^{\circ}_FG\to\mathbb{P}^1$ with general fibre $G$ and special fibre the normal cone $C_{F/G}$. Moreover, for any cartesian diagram
\begin{equation*}
\xymatrix{F^{\prime}\ar[r]^{f^{\prime}}\ar[d]&G^{\prime}\ar[d]\\
F\ar[r]^f&G}
\end{equation*}
there exists an induced morphism $M^{\circ}_{F^{\prime}}G^{\prime}\to M^{\circ}_FG$ that fits into a cartesian diagram
\begin{equation*}
\xymatrix{C_{F^{\prime}/G^{\prime}}\ar[r]\ar[d]&M^{\circ}_{F^{\prime}}G^{\prime}\ar[d]\\
C_{F/G}\ar[r]&M^{\circ}_FG}
\end{equation*}
\end{theorem}

\begin{remark}[{\cite[Remark 2.32]{manolache_VirtualPullbacks2012}}]\label{manolacheVirtualPullbacks2012-rem-2.32}
From Theorem \ref{manolacheVirtualPullbacks2012-thm-2.31} it follows that whenever $G$ is of pure dimension $r$, then $C_{F/G}$ is again of pure dimension $r$.
\end{remark}

\begin{proposition}[{\cite[Proposition 2.33, p.~14]{manolache_VirtualPullbacks2012}}]\label{manolacheVirtualPullbacks2012-prop-2.33}
\label{manolacheVirtualPullbacks2012-inc}
 If $f: F\to G$ is a DM-type morphism, then the cone stack $N_{F/G}$ of Definition \ref{manolacheVirtualPullbacks2012-def-2.30} is canonically isomorphic to the intrinsic normal sheaf $\mathfrak{N}_{F/G}$ of Definition \ref{manolacheVirtualPullbacks2012-def-2.12} and $C_{F/G}$ is canonically isomorphic to the intrinsic normal cone $\mathfrak{C}_{F/G}$ of Definition \ref{manolacheVirtualPullbacks2012-def-2.14}.
\end{proposition}

\begin{remark}[{\cite[Remark 2.34]{manolache_VirtualPullbacks2012}}]\label{manolacheVirtualPullbacks2012-rem-2.34}
 By the above Lemma we are allowed to identify the normal cone to a morphism with the intrinsic normal cone. In particular, the above Lemma shows that Definition \ref{manolacheVirtualPullbacks2012-def-2.30} is independent of the choice of $U$ and $V$ in diagram (11). Although normal cones are cone stacks, we will use for simplicity the notation $C_{F/G}$ instead of $\mathfrak{C}_{F/G}$.
\end{remark}

\begin{example}[{\cite[Example 2.35]{manolache_VirtualPullbacks2012}}]\label{manolacheVirtualPullbacks2012-ex-2.35}
Let $G$ be a DM-stack, $E$ a vector bundle on $G$ and $G\to E$ the zero section. Then $C_{G/E}$ is canonically isomorphic to $E$.
\end{example}

\begin{example}[{\cite[Example 2.36]{manolache_VirtualPullbacks2012}}]\label{manolacheVirtualPullbacks2012-ex-2.36}
Let $F\stackrel{f}\rightarrow G$ be a DM-type morphism and $p:E\to G$ a vector bundle on $G$. Let $i:G\to E$ be the zero section, and let $g:F\to E$ be $g:=i\circ f$. Then $C_{F/E}$ is canonically isomorphic to $C_{F/G}\times_Ff^*E$.
\end{example}

\begin{example}[{\cite[Example 2.37, p.~17]{manolache_VirtualPullbacks2012}}]\label{manolacheVirtualPullbacks2012-ex-2.37}
\label{manolacheVirtualPullbacks2012-vb}Let $F\stackrel{f}\rightarrow G$ be DM-type morphism, $\mathfrak{E}:=E^1/E^0$ a vector bundle stack on $G$. Let $G\stackrel{i}{\rightarrow}\mathfrak{E}$ denote the zero section. If $g:F\to G $ is the composition $F\stackrel{f}{\rightarrow}G\stackrel{i}{\rightarrow}\mathfrak{E}$, then $C_{F/G}$ is naturally isomorphic to $C_{F/G}\times_Ff^*\mathfrak{E}$
\end{example}

\begin{example}[{\cite[Example 2.38, p.~17]{manolache_VirtualPullbacks2012}}]\label{manolacheVirtualPullbacks2012-ex-2.38}
 \label{manolacheVirtualPullbacks2012-flat} Let $F\to G$ be a smooth morphism of DM-stacks. Then $C_{F/G}$ is isomorphic to $[F/T_{F/G}]$, hence it is a vector bundle stack.
\end{example}

\begin{example}[{\cite[Example 2.39]{manolache_VirtualPullbacks2012}}]\label{manolacheVirtualPullbacks2012-ex-2.39}
Let $X\stackrel{f}\rightarrow Y$ be a morphism of smooth schemes. Then, $U$ and $V$ above can be taken to be $X$ and $X\times Y$ as below
\begin{equation*}
\xymatrix{X\ar[r]^-{id\times f}\ar[d]&{X\times Y}\ar[d]^{\pi_2}\\
X\ar[r]&Y}
\end{equation*}
where $\pi_2$ is the projection on $Y$. It is then easy to see that the normal cone is $[N_{X/X\times Y}/T_{X}]$ that is a vector bundle stack.
\end{example}

\section{Construction}\label{manolacheVirtualPullbacks2012-sec-3}

\begin{proposition}[{\cite[Proposition 3.1, p.~18]{manolache_VirtualPullbacks2012}}]\label{manolacheVirtualPullbacks2012-prop-3.1}
(\cite{K:03}, Proposition 5.3.2)\label{manolacheVirtualPullbacks2012-stratif}
Let $F$ admit a stratification by global quotients  (see \cite{K:03}, Definition 3.5.3). Then, for any \emph{vector bundle stack} $\mathfrak{E}$, we have a canonical isomorphism $s^*:A_*(F)\to A_*(\mathfrak{E})$ .
\end{proposition}

\begin{remark}[{\cite[Remark 3.2]{manolache_VirtualPullbacks2012}}]\label{manolacheVirtualPullbacks2012-rem-3.2}
Every DM-stack admits a stratification by global quotients.
\\If $G$ admits a stratification by global quotients and $F\to G$ is a DM-type morphism then $F$ admits a stratification by global quotients.
\end{remark}

\subsection{Definition of virtual pull-backs}\label{manolacheVirtualPullbacks2012-sec-3.1}

\begin{condition}[{\cite[Condition 3.3]{manolache_VirtualPullbacks2012}}]\label{manolacheVirtualPullbacks2012-condition-3.3}
 We say that a morphism $f:F\to G$ of algebraic stacks and a vector bundle stack $\mathfrak{E}\to F$ satisfy condition ($\star$) if
\begin{enumerate}
\item $f$ is of DM-type,
\item we have fixed a closed embedding $i:C_{f}\hookrightarrow \mathfrak{E}$.
\end{enumerate}
\end{condition}

\begin{convention}[{\cite[Convention 3.4]{manolache_VirtualPullbacks2012}}]\label{manolacheVirtualPullbacks2012-convention-3.4}
Will say in short that the pair $(f,\mathfrak{E})$ satisfies condition ($\star$).
\end{convention}

\begin{remark}[{\cite[Remark 3.5]{manolache_VirtualPullbacks2012}}]\label{manolacheVirtualPullbacks2012-rem-3.5}
Let us consider a Cartesian diagram
\begin{equation*}
 \xymatrix {F^{\prime} \ar[r]\ar[d]_{p} & G^{\prime}\ar[d]^q\\
             F\ar[r]^f & G.}
\end{equation*}
If $\mathfrak{E}$ is a vector bundle on $F$ such that $C_{F/G}\hookrightarrow \mathfrak{E}$ is a closed embedding, then $C_{F^{\prime}/G^{\prime}}\hookrightarrow p^*\mathfrak{E}$ is a closed embedding.
\end{remark}

\begin{construct}[{\cite[Construction 3.6]{manolache_VirtualPullbacks2012}}]\label{manolacheVirtualPullbacks2012-constr-3.6}
Let $F$ be an Artin stack which admits a stratification by global quotient stacks and $\mathfrak{E}$ a vector bundle stack of (virtual) rank $n$ on $F$ such that $(f,\mathfrak{E})$ that satisfies condition $(\star)$ for $f$, we construct a pull-back map $f^{!}_\mathfrak{E}:A_*(G)\to A_{*-n}(F)$ as the composition
\begin{equation*}
 A_*(G)\stackrel{\sigma}\rightarrow A_*(C_{F/G})\stackrel{i_*}{\rightarrow}A_*(\mathfrak{E})\stackrel{s^*}{\rightarrow} A_{*-n}(F).
\end{equation*}
where
\begin{enumerate}
\item $\sigma$ is defined on the level of cycles by $\sigma(\sum n_i[V_i])=\sum n_i [C_{V_i\times_GF/V_i}]$
\item $i_*$ is the push-forward via the closed immersion $i$
\item $s^*$ is the morphism of Proposition \ref{manolacheVirtualPullbacks2012-prop-3.1}.
\end{enumerate}
By Proposition \ref{manolacheVirtualPullbacks2012-prop-2.26} we have a closed embedding of cones $$C_{V_i\times_GF/V_i}\hookrightarrow C_{F/G}.$$ The fact that $\sigma$ is well defined is a consequence of Theorem \ref{manolacheVirtualPullbacks2012-thm-2.31} (see \cite{K:03} Section 3 for local immersions and Section 5 for the general case).
 \\Going further, for any cartesian diagram
\begin{equation*}
 \xymatrix {F^{\prime} \ar[r]^{f^{\prime}}\ar[d]_{p} & G^{\prime}\ar[d]^q\\
             F\ar[r]^f & G}
\end{equation*}
such that $F^{\prime}$ admits a stratification by global quotient stacks and $\mathfrak{E}\to F$ satisfies condition $(\star)$ for $f$, let $f^{!}_{\mathfrak{E}}:A_*(G^{\prime})\to A_{*-n}(F^{\prime})$ be the composition
\begin{equation*}
 A_*(G^{\prime})\stackrel{\sigma}\rightarrow A_*(C_{F^{\prime}/G^{\prime}})\stackrel{\alpha}{\rightarrow} A_*(C_{F/G}\times_FF^{\prime})\stackrel{i_*}{\rightarrow}A_*(p^*\mathfrak{E})\stackrel{s^*}{\rightarrow}A_{*-n}(F^{\prime})
\end{equation*}
where $\alpha$ is the morphism from Proposition \ref{manolacheVirtualPullbacks2012-prop-2.26}.
\end{construct}

\begin{definition}[{\cite[Definition 3.7, p.~19]{manolache_VirtualPullbacks2012}}]\label{manolacheVirtualPullbacks2012-def-3.7}
\label{manolacheVirtualPullbacks2012-vp}
In the notation above, we call $f_{\mathfrak{E}}^!: A_*(G)\to A_*(F)$ a \textit{virtual pull-back}. When there is no risk of confusion we will omit the index.
\end{definition}

\begin{remark}[{\cite[Remark 3.8]{manolache_VirtualPullbacks2012}}]\label{manolacheVirtualPullbacks2012-rem-3.8}
In this remark we do not respect Convention \ref{manolacheVirtualPullbacks2012-convention-2.10}. If $E$ is a \emph{vector bundle} such that $(f,E)$ satisfies $(\star)$, then the above construction can be applied to any Artin stack $F$.  It is clear that in order to have $E$ a vector bundle $N_f$ must necessarily be a cone.
\\If $f$ is a locally closed embedding, then $N_f$ is a cone. Under this assumption if $E$ is a \emph{vector bundle} such that $(f,E)$ satisfies $(\star)$, then it is not necessary to ask $F$ to admit a stratification by global quotient stacks (see \cite{K:02}, Theorem 2.1.12 (vi)).
\end{remark}

\begin{remark}[{\cite[Remark 3.9]{manolache_VirtualPullbacks2012}}]\label{manolacheVirtualPullbacks2012-rem-3.9}
Note that in case $X$, $Y$ are schemes such that $X$ is regularly embedded in $Y$, then the normal bundle of $X$ in $Y$ satisfies condition $\star$ and $i^!_{N_{X/Y}}$ is precisely the refined Gysin pull-back of \cite{f}, Section 6.2. We remark that the pull-back depends on the chosen bundle. For example, if $(i,E)$ satisfies condition $(\star)$ we can construct $i^!_{E\oplus E^{\prime}}$, where $E^{\prime}$ is any other vector bundle. These morphisms will be obviously different from each other.
\end{remark}

\begin{remark}[{\cite[Remark 3.10, p.~20]{manolache_VirtualPullbacks2012}}]\label{manolacheVirtualPullbacks2012-rem-3.10}
\label{manolacheVirtualPullbacks2012-fultflat}
 If $f:X\to Y$ is a smooth morphism of schemes, then by Example \ref{manolacheVirtualPullbacks2012-ex-2.38} $C_{X/Y}$ is a vector bundle stack and hence we can construct the  associated virtual pull-back $f^!_{C_{X/Y}}:A_*(Y)\to A_*(X)$. This has already been defined in \cite{K:03} and it is proved that the definition agrees with the usual flat pull-back (see e.g. \cite{f}).
\end{remark}

\begin{proposition}[{\cite[Proposition 3.11]{manolache_VirtualPullbacks2012}}]\label{manolacheVirtualPullbacks2012-prop-3.11}
\uses{manolacheVirtualPullbacks2012-prop-2.33}
If $F\stackrel{f}\rightarrow G$ is a DM-type morphism and there exists a perfect relative obstruction theory $E_{F/G}^{\bullet}$, then condition $(\star)$ is fulfilled.
\\Conversely, if $F\stackrel{f}\rightarrow G$ is a morphism that satisfies condition ($\star$), then there exists a perfect obstruction theory $E^{\bullet}_{F/G}\to L_{F/G}$ such that $\mathfrak{E}=h^1/h^0({E^{\bullet}_{F/G}}^{\vee})$ which is unique up to quasi-isomorphism.
\end{proposition}

\begin{corollary}[{\cite[Corollary 3.12, p.~20]{manolache_VirtualPullbacks2012}}]\label{manolacheVirtualPullbacks2012-cor-3.12}
\label{manolacheVirtualPullbacks2012-extrasmooth}
If $F\stackrel{f}\rightarrow G$ is a DM-type morphism such that there exists a perfect relative obstruction theory $E^{\bullet}_{F/G}$ and $G$ is a stack of pure dimension, then $f^!_{E^{\bullet}_{F/G}}([G])$ is a virtual class of $F$ in the sense of \cite{bf}.
\end{corollary}

\subsection{A fundamental example of Obstruction Theory}\label{manolacheVirtualPullbacks2012-sec-3.2}

\begin{construct}[{\cite[Construction 3.13, p.~21]{manolache_VirtualPullbacks2012}}]\label{manolacheVirtualPullbacks2012-constr-3.13}
\label{manolacheVirtualPullbacks2012-constr}
 Let $f:F\to G$ be a DM-type morphism and let $F$ and $G$ be DM-stacks having relative obstruction theories with respect to some smooth Artin stack $\mathfrak{M}$. Let us denote them by $E_{F/\mathfrak{M}}^{\bullet}$ and $E_{G/\mathfrak{M}}^{\bullet}$ respectively. Given a morphism $\varphi: f^*E_{G/\mathfrak{M}}^{\bullet}\to E_{F/\mathfrak{M}}^{\bullet}$ commuting with $f^*L_{G/\mathfrak{M}}\to L_{F/\mathfrak{M}}$, we construct a relative obstruction theory $E_{F/G}^{\bullet}$.
\\The morphism $f:F\to G$ induces a distinguished triangle of cotangent complexes
\begin{equation*}
 f^*L_{G/\mathfrak{M}}\to L_{F/\mathfrak{M}}\to L_{F/G}\to f^*L_{G/\mathfrak{M}}[1].
\end{equation*}
Similarly, $\varphi$ gives rise to a distinguished triangle
\begin{equation}\label{manolacheVirtualPullbacks2012-obstr}
 f^*E_{G/\mathfrak{M}}^{\bullet}\stackrel{\varphi}{\rightarrow} E_{F/\mathfrak{M}}^{\bullet}\to E_{F/G}^{\bullet}\to f^*E_{G/\mathfrak{M}}[1]
\end{equation}

hence we have a morphism of distinguished triangles that induces the following morphism in cohomology
\begin{equation*}
\scriptsize
\xymatrix@C=0.43cm{h^{-1}(f^*E_{G/\mathfrak{M}}^{\bullet}) \ar[r]\ar[d] & h^{-1}(E_{F/\mathfrak{M}}^{\bullet})\ar[d]\ar[r]\ar[d] & h^{-1}(E_{F/G}^{\bullet})\ar[r]\ar[d] & h^0(f^*E_{G/\mathfrak{M}}^{\bullet}) \ar[r]\ar[d] & h^0(E_{F/\mathfrak{M}}^{\bullet})\ar[d]\ar[r]\ar[d] & h^0(E_{F/G}^{\bullet})\ar[d]\\h^{-1}(f^*L_{G/\mathfrak{M}}^{\bullet}) \ar[r] & h^{-1}(L_{F/\mathfrak{M}}^{\bullet})\ar[r] & h^{-1}(L_{F/G}^{\bullet})\ar[r]& h^0(f^*L_{G/\mathfrak{M}}^{\bullet}) \ar[r] & h^0(L_{F/\mathfrak{M}}^{\bullet})\ar[r] & h^0(L_{F/G}^{\bullet})}
\end{equation*}
We know that the first two vertical arrows are surjective and by the definition of obstruction theories we get by a simple diagram chase that $E_{F/G}^{\bullet}$ is also an obstruction theory.
\end{construct}

\begin{remark}[{\cite[Remark 3.14]{manolache_VirtualPullbacks2012}}]\label{manolacheVirtualPullbacks2012-rem-3.14}
 Let us note that the morphism of mapping cones $E_{F/G}^{\bullet}\to L_{F/G}$ is not unique in $\mathcal{D}^{\leq0}_F$ and therefore this procedure gives many \emph{different} relative obstruction theories.
\end{remark}

\begin{remark}[{\cite[Remark 3.15, p.~21]{manolache_VirtualPullbacks2012}}]\label{manolacheVirtualPullbacks2012-rem-3.15}
\label{manolacheVirtualPullbacks2012-smooth}
 If $G$ is smooth over $\mathfrak{M}$ and $E_{G/\mathfrak{M}}^{\bullet}$ is trivial (i.e. $E_{G/\mathfrak{M}}^{\bullet}=L_{G/\mathfrak{M}}^{\bullet}$), then the above diagram shows that $E_{F/G}$ is perfect in $[-1,0]$.
\end{remark}

\begin{example}[{\cite[Example 3.16, p.~21]{manolache_VirtualPullbacks2012}}]\label{manolacheVirtualPullbacks2012-ex-3.16}
\label{manolacheVirtualPullbacks2012-easyfunct} A special case of this construction is when $F\to G$ is a locally closed immersion and $G$ is taken to be smooth over $\mathfrak{M}$. Taking $h^{-1}(f^*E_{G/\mathfrak{M}}^{\bullet})=0$ we obtain that $h^{-2}(E_{F/G}^{\bullet})=0$. This makes $E_{F/G}^{\bullet}$ into a perfect obstruction theory concentrated in degree $-1$ and consequently $\mathfrak{E}$ into a vector bundle.
\end{example}

\begin{example}[{\cite[Example 3.17]{manolache_VirtualPullbacks2012}}]\label{manolacheVirtualPullbacks2012-ex-3.17}
 The basic case.
\\In the notation above, let us suppose $G$ is smooth and $F\stackrel{i}{\hookrightarrow}G$ is a closed substack and there exists a morphism $f^*L_G\stackrel{\varphi}{\rightarrow}E_{F}^{\bullet}$. Suppose $\mathfrak{M}$ has pure dimension. Let $E_{F/G}^{\bullet}$ be the cone of $\varphi$. By Construction \ref{manolacheVirtualPullbacks2012-constr-3.13} it is a perfect obstruction theory for $i$. Then we have
\\(i) $(C_{F/G},E_{F/G}^{\bullet})$ induces the same virtual class on $F$ as $(C_{F/\mathfrak{M}},E_{F/\mathfrak{M}}^{\bullet})$.
\\(ii) The pull back defined by $E_{F/G}^{\bullet}$ respects the relation $$i^![G]=[F]^{\mathrm{virt}}.$$
\end{example}

\section{Basic properties}\label{manolacheVirtualPullbacks2012-sec-4}

\begin{theorem}[{\cite[Theorem 4.1, p.~23]{manolache_VirtualPullbacks2012}}]\label{manolacheVirtualPullbacks2012-thm-4.1}
\uses{manolacheVirtualPullbacks2012-prop-2.26}
 \label{manolacheVirtualPullbacks2012-relations}Consider a fibre diagram of Artin stacks
\begin{equation*}
 \xymatrix {F^{\prime\prime}\ar[r]\ar[d]_{q} & G^{\prime\prime}\ar[d]^p\\
F^{\prime} \ar[r]^{f^{\prime}}\ar[d]_{g} & G^{\prime}\ar[d]^h\\
             F\ar[r]^f & G}
\end{equation*}
and let us assume
\begin{enumerate}
\item $F^{\prime}$, $F^{\prime\prime}$ admit stratifications by global quotient stacks,
\item ${\mathfrak{E}}$ is a vector bundle stack of rank $d$ such that $(f,{\mathfrak{E}})$ satisfies condition $(\star)$ for $f$.
\item $\mathfrak{E}$ is isomorphic to a global quotient $[E^1/E^0]$, with $E^1$, $E^0$ vector bundles on $F$.

\end{enumerate}
(i) (Push-forward) If $p$ is either a projective morphism of Artin stacks or a proper morphism of DM-stacks and $\alpha\in A_k(G^{\prime\prime})$, then $f^!_{\mathfrak{E}}p_*(\alpha)=q_*f^!_{\mathfrak{E}}\alpha$ in $A_{k-d}(F^{\prime})$.\\
(ii)\ (Pull-back) If $p$ is flat of relative dimension $n$ and $\alpha\in A_k(G^{\prime})$, then $f^!_{\mathfrak{E}}p^*(\alpha)=q^*f^!_{\mathfrak{E}}\alpha$ in $A_{k+n-d}(F^{\prime\prime})$\\
(iii)(Compatibility) If $\alpha\in A_k(G^{\prime\prime})$, then $f^{!}_{\mathfrak{E}}\alpha=f^{\prime!}_{g^*{\mathfrak{E}}}\alpha$ in $A_{k-d}(F^{\prime\prime})$.
\end{theorem}

\begin{remark}[{\cite[Remark 4.2]{manolache_VirtualPullbacks2012}}]\label{manolacheVirtualPullbacks2012-rem-4.2}
 If $p$ is projective we do not need $\mathfrak{E}$ to be a global quotient. The complication in Step 2 of the proof of Theorem \ref{manolacheVirtualPullbacks2012-thm-4.1} (i) is due to the fact that push-forwards along proper morphisms  of \emph{Artin stacks} cannot be defined unless the morphism is projective. Concretely, if $q$ is proper but not projective, we do not know how to define $Q_*$ in diagram (19).
\\On the contrary, if $f$ is a local embedding of DM stacks and $q$ is proper then, the deformation spaces are DM stacks and $P$ is proper; this implies that we can push-forward cycles along $P$ and $Q$.
\end{remark}

\begin{theorem}[{\cite[Theorem 4.3]{manolache_VirtualPullbacks2012}}]\label{manolacheVirtualPullbacks2012-thm-4.3}
\uses{manolacheVirtualPullbacks2012-thm-4.1}
(Commutativity)
Consider a fiber diagram of Artin stacks
\begin{equation*}
\xymatrix{F^{\prime\prime}\ar[d]_q\ar[r]^v&G^{\prime\prime}\ar[d]\ar[r]^u&H\ar[d]^g\\
F^{\prime}\ar[d]_p\ar[r]&G{^\prime}\ar[d]\ar[r]&K\\
F\ar[r]^f&G}
\end{equation*}such that $F^{\prime}$ and $G^{\prime\prime}$ admit stratifications by global quotients. Let us assume $f$ and $g$ are morphisms of DM-type and let ${\mathfrak{E}}$ and ${\mathfrak{F}}$ be vector bundle stacks of rank $d$ and $e$ respectively such that $(f,{\mathfrak{E}})$ and $(g,{\mathfrak{F}})$ satisfy condition $(\star)$. Then for all $\alpha\in A_k(G^{\prime})$,
\begin{equation*}
 g^!_{\mathfrak{F}}f^!_{\mathfrak{E}}(\alpha)= f^!_{\mathfrak{E}}g^!_{\mathfrak{F}}(\alpha)
\end{equation*}
in $A_{k-d-e}(F^{\prime\prime})$.
\end{theorem}

\begin{theorem}[{\cite[Theorem 4.4, p.~27]{manolache_VirtualPullbacks2012}}]\label{manolacheVirtualPullbacks2012-thm-4.4}
\label{manolacheVirtualPullbacks2012-bivar}
Let $F$ be an Artin stacks which admits a stratification by global quotients, let $f:F\to G$ be a morphism and $\mathfrak{E}\to F$ be a rank-$n$ vector bundle stack on $F$ such that $(f,\mathfrak{E})$ satisfies  Condition ($\star$). Then $f_{\mathfrak{E}}^!$ defines a bivariant class in $A^{n}(F\to G)$ in the sense of \cite{f}, Definition 17.1.
\end{theorem}

\begin{definition}[{\cite[Definition 4.5, p.~27]{manolache_VirtualPullbacks2012}}]\label{manolacheVirtualPullbacks2012-def-4.5}
\label{manolacheVirtualPullbacks2012-compatib}
 Let $F\stackrel{f}\rightarrow G\stackrel{g}\rightarrow \mathfrak{M}$ be DM-type morphisms of stacks. If we are given a distinguished triangle of relative obstruction theories which are perfect in $[-1,0]$ $$g^*E^{\bullet}_{G/\mathfrak{M}}\stackrel{\varphi}\rightarrow E^{\bullet}_{F/\mathfrak{M}}\to E^{\bullet}_{F/G}\to g^*E^{\bullet}_{G/\mathfrak{M}}[1]$$ with a morphism to the distinguished triangle $$g^*L_{G/\mathfrak{M}}\to L_{F/\mathfrak{M}}\to L_{F/G}\to g^*L_{G/\mathfrak{M}}[1],$$ then we call $(E^{\bullet}_{F/G},E^{\bullet}_{G/\mathfrak{M}},E^{\bullet}_{F/\mathfrak{M}})$ a compatible triple.
\end{definition}

\begin{remark}[{\cite[Remark 4.6]{manolache_VirtualPullbacks2012}}]\label{manolacheVirtualPullbacks2012-rem-4.6}
 As in Construction \ref{manolacheVirtualPullbacks2012-constr-3.13}, if there is a morphism $E^{\bullet}_{F/G}\stackrel{\psi}\rightarrow g^*E^{\bullet}_{G/\mathfrak{M}}[1]$ compatible with the corresponding morphism between the cotangent complexes, then $\psi$ determines a complex $E^{\bullet}_{F/\mathfrak{M}}$ which fits in a distinguished triangle as above. Moreover, $E^{\bullet}_{F/\mathfrak{M}}$ defines a relative obstruction theory. If $E^{\bullet}_{F/G}$ and $E^{\bullet}_{G/\mathfrak{M}}$ are perfect, then $E^{\bullet}_{F/\mathfrak{M}}$ is perfect.
\end{remark}

\begin{lemma}[{\cite[Lemma 4.7]{manolache_VirtualPullbacks2012}}]\label{manolacheVirtualPullbacks2012-lem-4.7}
\uses{manolacheVirtualPullbacks2012-ex-2.37, manolacheVirtualPullbacks2012-thm-4.1}
Consider a fibre diagram
\begin{equation*}
 \xymatrix {F^{\prime} \ar[r]^{f^{\prime}}\ar[d]_{p} & G^{\prime}\ar[d]^q\ar[r]^{0}&\mathfrak{F}^{\prime}\ar[d]^r\\
             F\ar[r]^f & G\ar[r]^0&\mathfrak{F}.}
\end{equation*}
with $F^{\prime}$ an Artin stack which admits a stratification by global quotients, $\pi:\mathfrak{F}\to G$ a vector bundle stack of rank $e$ on $G$ and $\pi^{\prime}:\mathfrak{F}^{\prime}\to G^{\prime}$ its pullback to $G^{\prime}$. Let us assume $\mathfrak{E}^{\prime}$ is a vector bundle stack on $F$ such that $(f,\mathfrak{E}^{\prime})$ satisfies condition $(\star)$. Then we have a natural map $$C_{F/\mathfrak{F}}\to\mathfrak{E}:=\mathfrak{E}^{\prime}\oplus f^*\mathfrak{F}$$ which is a closed embedding and for any $\alpha\in A_k(\mathfrak{F}^{\prime})$ $$(0\circ f)_{\mathfrak{E}}^!(\alpha)=f^!_{\mathfrak{E}^{\prime}}(0^!_{\mathfrak{F}}(\alpha)).$$
\end{lemma}

\begin{theorem}[{\cite[Theorem 4.8, p.~29]{manolache_VirtualPullbacks2012}}]\label{manolacheVirtualPullbacks2012-thm-4.8}
\uses{manolacheVirtualPullbacks2012-constr-3.13, manolacheVirtualPullbacks2012-def-4.5}
\label{manolacheVirtualPullbacks2012-functoriality}
(Functoriality)
 Consider a fibre diagram
\begin{equation*}
 \xymatrix {F^{\prime} \ar[r]^{f^{\prime}}\ar[d]_{p} & G^{\prime}\ar[d]^q\ar[r]^{g^{\prime}}&\mathfrak{M}^{\prime}\ar[d]^r\\
             F\ar[r]^f & G\ar[r]^g&\mathfrak{M}.}
\end{equation*}
Let us assume $f$, $g$ and $g\circ f$ are DM-type morphisms and have perfect relative obstruction theories $E^{\prime\bullet}$, $E^{\prime\prime\bullet}$ and $E^{\bullet}$ respectively and let us denote the associated vector bundle stacks by $\mathfrak{E}^{\prime}$, $\mathfrak{E}^{\prime\prime}$ and $\mathfrak{E}$ respectively. If $F^{\prime},\ F^{\prime}$ admit stratifications by global quotients and $(E^{\prime\bullet}, E^{\prime\prime\bullet},E^{\bullet})$ is a compatible triple, then for any $\alpha\in A_k(\mathfrak{M}^{\prime})$ $$(g\circ f)_{\mathfrak{E}}^!(\alpha)=f^!_{\mathfrak{E}^{\prime}}(g^!_{\mathfrak{E}^{\prime\prime}}(\alpha)).$$
\end{theorem}

\begin{corollary}[{\cite[Corollary 4.9, p.~31]{manolache_VirtualPullbacks2012}}]\label{manolacheVirtualPullbacks2012-cor-4.9}
\uses{manolacheVirtualPullbacks2012-thm-4.8}
\label{manolacheVirtualPullbacks2012-virtual}
Let us assume we have a commutative diagram
\begin{equation*}
 \xymatrix{F\ar[rr]^f\ar[rd]_{\epsilon_F}&&G\ar[ld]^{\epsilon_{G}}\\
&\mathfrak{M}}
\end{equation*}
with $\mathfrak{M}$ of pure dimension. If $F$ and $G$ admit stratifications by global quotients and  $(E^{\bullet}_{F/G},E^{\bullet}_{G/\mathfrak{M}},E^{\bullet}_{F/\mathfrak{M}})$ a compatible triple, then $$f_{\mathfrak{E}_{F/G}}^![G]^{\mathrm{virt}}=[F]^{\mathrm{virt}}.$$
\end{corollary}

\begin{remark}[{\cite[Remark 4.10]{manolache_VirtualPullbacks2012}}]\label{manolacheVirtualPullbacks2012-rem-4.10}
Let us consider a cartesian diagram of DM stacks
\begin{equation*}
 \xymatrix {F^{\prime} \ar[r]\ar[d]_{g} & G^{\prime}\ar[d]^f\\
             F\ar[r]^i & G}
\end{equation*}
with obstruction theories $E_{F}^{\bullet}$, $E_{G}^{\bullet}$, $E_{F^{\prime}}^{\bullet}$, $E_{G^{\prime}}^{\bullet}$. Let us assume we have perfect obstruction theories $E_{F/G}^{\bullet}$ and $E_{F^{\prime}/G^{\prime}}^{\bullet}$ compatible with  $E_{F}^{\bullet}$, $E_{G}^{\bullet}$ and $E_{F^{\prime}}^{\bullet}$, $E_{G^{\prime}}^{\bullet}$. If $g^*\mathfrak{E}_{F/G}=\mathfrak{E}_{F^{\prime}/G^{\prime}}$, then $i^![G^{\prime}]^{\mathrm{virt}}=[F^{\prime}]^{\mathrm{virt}}$.
\\This generalizes Proposition 5.10 in \cite{bf} and Theorem 1 in \cite{K:01}.
\end{remark}

\section{Applications}\label{manolacheVirtualPullbacks2012-sec-5}

\subsection{Preliminaries}\label{manolacheVirtualPullbacks2012-sec-5.1}

\begin{remark}[{\cite[Remark 5.1]{manolache_VirtualPullbacks2012}}]\label{manolacheVirtualPullbacks2012-rem-5.1}
Let $f: X\to Y$ be a morphism of smooth algebraic varieties. Let $\beta\in A_1(X)$ and $g,\ n$ be any natural numbers. Then $f$ induces a morphism of stacks $\bar{f}: \overline{M}_{g,n}(X,\beta)\to \overline{M}_{g,n}(Y,f_*\beta)$.
\\Convention: Given a morphism of smooth algebraic varieties $f:X\to Y$, we will indicate the induced morphism between moduli spaces of stable maps by the same letter with a bar.
\end{remark}

\begin{theorem}[{\cite[Theorem 5.2, p.~33]{manolache_VirtualPullbacks2012}}]\label{manolacheVirtualPullbacks2012-thm-5.2}
 (Cohomology and Base Change)\label{manolacheVirtualPullbacks2012-grothmaps} Let $G^{\prime}\stackrel{p}{\rightarrow}G$ be a flat morphism of separated DM stacks and $E^{\bullet}$ a finite bounded of locally free sheaf on $G^{\prime}$. Then for any base change
\begin{equation*}
 \xymatrix{F^{\prime}\ar[r]^{g}\ar[d]_q&G^{\prime}\ar[d]^p\\
F\ar[r]^f&G}
\end{equation*}
we have a canonical isomorphism in $\mathcal{D}_F$ $$Lf^*Rp_*E^{\bullet}\simeq Rq_*Lg^*E^{\bullet}.$$
\end{theorem}

\begin{proposition}[{\cite[Proposition 5.3, p.~33]{manolache_VirtualPullbacks2012}}]\label{manolacheVirtualPullbacks2012-prop-5.3}
\uses{manolacheVirtualPullbacks2012-thm-5.2, manolacheVirtualPullbacks2012-rem-3.15}
\label{manolacheVirtualPullbacks2012-cbg}
 Let $f:X\rightarrow \mathbb{P}$ be a morphism of smooth projective varieties and let $T_{X/\mathbb{P}}$ be the dual of the cotangent complex of $X$ to $\mathbb{P}$. Then, in notations as above
\\(i) The map $\bar{f}:\overline{M}_{g,n}(X,\beta)\to \overline{M}_{g,n}(\mathbb{P},f_*\beta)$ has a dual obstruction theory $E_{\overline{M}_{g,n}(X,\beta)/\overline{M}_{g,n}(\mathbb{P},f_*\beta)}^{\bullet}$ isomorphic to $\mathcal{R}^{\bullet}{\pi_X}_*ev_X^*T_{X/\mathbb{P}}$ in $\mathcal{D}_{\overline{M}_{g,n}(X,\beta)}$.
\\(ii) If $f=i$ is an embedding and $N_{X/\mathbb{P}}$ denotes the normal bundle of $X$ in $\mathbb{P}$, then $\bar{i}:\overline{M}_{g,n}(X,\beta)\to \overline{M}_{g,n}(\mathbb{P},f_*\beta)$ has a dual obstruction theory $(E_{\overline{M}_{g,n}(X,\beta)/\overline{M}_{g,n}(\mathbb{P},i*\beta)}^{\bullet})^{\vee}$ $$[0\to\mathcal{R}^0{\pi_X}_*ev_X^* N_{X/\mathbb{P}}\to \mathcal{R}^1{\pi_X}_*ev_X^* N_{X/\mathbb{P}}]$$ in $[0,2]$.
\\(iii) If $g=0$ and $\mathbb{P}$ is convex, then $E_{\overline{M}_{g,n}(X,\beta)/\overline{M}_{g,n}(\mathbb{P},f_*\beta)}^{\bullet}$ is perfect.
\end{proposition}

\begin{proposition}[{\cite[Proposition 5.4, p.~34]{manolache_VirtualPullbacks2012}}]\label{manolacheVirtualPullbacks2012-prop-5.4}
\uses{manolacheVirtualPullbacks2012-prop-5.3, manolacheVirtualPullbacks2012-thm-5.2}
\label{manolacheVirtualPullbacks2012-cartesian}
Let
\begin{equation}\label{manolacheVirtualPullbacks2012-hopefullycart}
 \xymatrix{X^{\prime}\ar[r]^j\ar[d]^p&\mathbb{P}^{\prime}\ar[d]^q
\\X\ar[r]^i&\mathbb{P}}
\end{equation}
be a cartesian diagram of smooth projective varieties and let $\beta\in A_1(X^{\prime})$ be any homology class of a curve.
\\(i) Then the induced diagram of moduli spaces of stable maps
\begin{equation*}
 \xymatrix{\overline{M}_{g,n}(X^{\prime},\beta)\ar[r]^{\bar{j}}\ar[d]^{\bar{p}}&\overline{M}_{g,n}(\mathbb{P}^{\prime},j_*\beta)\ar[d]^{\bar{q}}
\\\overline{M}_{g,n}(X,p_*\beta)\ar[r]^{\bar{i}}&\overline{M}_{g,n}(\mathbb{P}, (i\circ p)_*\beta)}
\end{equation*}
is commutative. If $i$ is a closed embedding, then it induces an open and closed embedding of $\overline{M}_{g,n}(X^{\prime},\beta)$ in the fiber product $\overline{M}_{g,n}(X,p_*\beta)\times_{\overline{M}_{g,n}(\mathbb{P}, (i\circ p)_*\beta)}\overline{M}_{g,n}(\mathbb{P}^{\prime},j_*\beta)$.
\\(ii)If the natural map $p^*L_{X/\mathbb{P}}\to L_{X^{\prime}/\mathbb{P}^{\prime}}$ is an isomorphism, then it induces an isomorphism $$E_{\overline{M}_{g,n}(X^{\prime},\beta)/\overline{M}_{g,n}(\mathbb{P}^{\prime},j_*\beta)}^{\bullet}\simeq\bar{p}^*E_{\overline{M}_{g,n}(X,p_*\beta)/\overline{M}_{g,n}(\mathbb{P},(i\circ p)_*\beta)}^{\bullet}.$$ If moreover, $g=0$ and $\mathbb{P}$ is convex then, $E_{\overline{M}_{g,n}(X^{\prime},\beta)/\overline{M}_{g,n}(\mathbb{P}^{\prime},j_*\beta)}^{\bullet}$ is perfect.
\end{proposition}

\begin{example}[{\cite[Example 5.5]{manolache_VirtualPullbacks2012}}]\label{manolacheVirtualPullbacks2012-ex-5.5}
Let us consider a cartesian diagram of smooth projective varieties as above. The induced commutative diagram of moduli spaces of stable maps will not be cartesian in general. One counterexample is the case of $\mathbb{P}$ is a point $P:=Spec\ k$, $\mathbb{P}^{\prime}:=Y$ and $X^{\prime}:=X\times Y$. The diagram
\begin{equation*}
 \xymatrix{\overline{M}_{g,n}(X\times Y,(\beta_1, \beta_2))\ar[r]^-{\bar{p_2}}\ar[d]^{\bar{p_1}}&\overline{M}_{g,n}(Y,\beta_2)\ar[d]
\\\overline{M}_{g,n}(X,\beta_1)\ar[r]&\overline{M}_{g,n}(P, 0)}
\end{equation*}
is \emph{not} cartesian. Let $M$ be the cartesian product $\overline{M}_{g,n}(Y,\beta_2)\times_{\overline{M}_{g,n}(P, 0)}\overline{M}_{g,n}(X,\beta_1)$. Then, we still have a nice relation between the virtual class of $M$ and the virtual class of $\overline{M}_{g,n}(X\times Y,(\beta_1, \beta_2))$. This was studied by Behrend (see \cite{b2}, Theorem 1).
\end{example}

\begin{example}[{\cite[Example 5.6]{manolache_VirtualPullbacks2012}}]\label{manolacheVirtualPullbacks2012-ex-5.6}
Let
\begin{equation*}
 \xymatrix{X^{\prime}\ar[r]^j\ar[d]^p&\mathbb{P}^{\prime}\ar[d]^q
\\X\ar[r]^i&\mathbb{P}}
\end{equation*}
be a cartesian diagram of smooth projective varieties as in the above proposition. Let us moreover suppose that $i$ and $j$ induce injective morphisms of groups $i_*:A_1(X)\to A_1(\mathbb{P})$, respectively $j_*:A_1(X^{\prime})\to A_1(\mathbb{P}^{\prime})$ (but $i$ is not supposed to be injective!). Then, the corresponding commutative diagram between moduli spaces of stable maps is cartesian.
\end{example}

\begin{remark}[{\cite[Remark 5.7, p.~36]{manolache_VirtualPullbacks2012}}]\label{manolacheVirtualPullbacks2012-rem-5.7}
\label{manolacheVirtualPullbacks2012-condand}
 If $i$ is a closed embedding and the natural map $A_1(X^{\prime})\to A_1(\mathbb{P}^{\prime})\oplus A_1(X)$ is injective then diagram (\ref{manolacheVirtualPullbacks2012-hopefullycart}) is cartesian. This follows from the fact that the above condition implies that the natural map  $$\overline{M}_{g,n}(X^{\prime},\beta)\to\overline{M}_{g,n}(X,p_*\beta)\times_{\overline{M}_{g,n}(\mathbb{P}, (i\circ p)_*\beta)}\overline{M}_{g,n}(\mathbb{P}^{\prime},j_*\beta)$$ is surjective.
\end{remark}

\begin{example}[{\cite[Example 5.8]{manolache_VirtualPullbacks2012}}]\label{manolacheVirtualPullbacks2012-ex-5.8}
Let us now look at an example where our construction of virtual pull-backs does not apply. If we consider $\overline{M}_{g,n}(X,\beta)\stackrel{\bar{i}}{\rightarrow}\overline{M}_{g,n}(\mathbb{P}, i_*\beta)$ with $\mathbb{P}$ a convex space, then the construction applies without further conditions only in genus zero. In general (higher genus or non-convex $\mathbb{P}$) $h^{-2}(E_{\overline{M}_{g,n}(X,\beta)/\overline{M}_{g,n}(\mathbb{P}, i_*\beta)})$ might not vanish.
\\To see an example, let us consider $i:\mathbb{P}^r\hookrightarrow\mathbb{P}^r\times\mathbb{P}^s$, the inclusion into the first factor. Then we have an induced map $\overline{M}_{g,n}(\mathbb{P}^r,d_1)\hookrightarrow\overline{M}_{g,n}(\mathbb{P}^r\times\mathbb{P}^s,(d_1,0))$. From Corollary 5.3 we obtain that the dual relative obstruction theory of $\\bar{i}$ is $\mathcal{R}^{\bullet}\pi_*ev^*N_{\mathbb{P}^r/\mathbb{P}^r\times\mathbb{P}^s}$. We have that the normal bundle $N_{\mathbb{P}^r/\mathbb{P}^r\times\mathbb{P}^s}$ is isomorphic to $\mathcal{O}_{\mathbb{P}^r}^{\oplus s}$. Since $\mathcal{R}^{1}\pi_*f^*(\mathcal{O}_{\mathbb{P}^r}^{\oplus s})$ is non-zero for $g\geq1$, the (dual) relative obstruction theory will never be perfect.
\end{example}

\subsection{Pulling back divisors}\label{manolacheVirtualPullbacks2012-sec-5.2}

\begin{remark}[{\cite[Remark 5.9]{manolache_VirtualPullbacks2012}}]\label{manolacheVirtualPullbacks2012-rem-5.9}
One could na\"ively hope to obtain new relations between the rational GW invariants of $X$ by pulling back the WDVV relations in $\mathbb{P}$. The above shows that for any $X\hookrightarrow\mathbb{P}$ pulling-back the WDVV equations in $\overline{M}_{0,n}(\mathbb{P},d)$ gives the WDVV equations in $\overline{M}_{0,n}(X,d)$.
\end{remark}

\subsection{Blow-ups}\label{manolacheVirtualPullbacks2012-sec-5.3}

\begin{definition}[{\cite[Definition 5.10, p.~38]{manolache_VirtualPullbacks2012}}]\label{manolacheVirtualPullbacks2012-def-5.10}
 \label{manolacheVirtualPullbacks2012-lifting}
For a blow up $p_X:\tilde{X}\to X$ and a class $\beta\in A_1(X)$ we call the class $p_X^!\beta$ the lifting of $\beta$ and we denote it by $\tilde{\beta}$, where $p_X^!$ is the refined intersection product of \cite{f}, Chapter 8.
\end{definition}

\begin{remark}[{\cite[Remark 5.11]{manolache_VirtualPullbacks2012}}]\label{manolacheVirtualPullbacks2012-rem-5.11}
 The lifting of $\beta$ satisfies two basic properties that follow trivially from the projection formula, namely $(p_X)_*\tilde{\beta}=\beta$ and $\tilde{\beta}\cdot E=0$.
\end{remark}

\begin{lemma}[{\cite[Lemma 5.12]{manolache_VirtualPullbacks2012}}]\label{manolacheVirtualPullbacks2012-lem-5.12}
The moduli space of stable maps to $\tilde{X}$ of class $\tilde{\beta}$ and the moduli space of stable maps to $X$ of class $\beta$ have the same virtual dimension.
\end{lemma}

\begin{lemma}[{\cite[Lemma 5.13, p.~38]{manolache_VirtualPullbacks2012}}]\label{manolacheVirtualPullbacks2012-lem-5.13}
\label{manolacheVirtualPullbacks2012-andreas}
 In notations as before, the natural projection $p_{X}:\tilde{X}\to X$ induces a morphism $\bar{p}_{X}:\overline{M}_{0,n}(\tilde{X},\tilde{\beta})\to\overline{M}_{0,n}(X,\beta)$. If $X=\mathbb{P}$ is a homogeneous space and $\tilde{X}=\tilde{\mathbb{P}}$ the blow up of $\mathbb{P}$, then $$({\bar{p}}_{\mathbb{P}})_*[\overline{M}_{0,n}(\tilde{\mathbb{P}},\tilde{\beta})]^{\mathrm{virt}}=[\overline{M}_{0,n}(\mathbb{P},\beta)]^{\mathrm{virt}}.$$
\end{lemma}

\begin{proposition}[{\cite[Proposition 5.14, p.~39]{manolache_VirtualPullbacks2012}}]\label{manolacheVirtualPullbacks2012-prop-5.14}
\uses{manolacheVirtualPullbacks2012-thm-4.1, manolacheVirtualPullbacks2012-rem-5.7, manolacheVirtualPullbacks2012-constr-3.13, manolacheVirtualPullbacks2012-cor-4.9, manolacheVirtualPullbacks2012-prop-5.4, manolacheVirtualPullbacks2012-lem-5.13}
\label{manolacheVirtualPullbacks2012-main} Let $X$ be a smooth projective subvariety of some homogeneous space $\mathbb{P}$ and let us assume that there exists $Z$ a smooth subvariety of $\mathbb{P}$, such that $X$ and $Z$ intersect transversely. Let $Y:=X\cap Z$ and $\tilde{X}$ be the blow-up of $X$ along $Y$. Then for any non-negative integer $n$ and any $\beta\in A_1(X)$ with lifting $\tilde{\beta}\in A_*(\tilde{X})$ $$(\bar{p}_{X})_*[\overline{M}_{0,n}(\tilde{X},\tilde{\beta})]^{\mathrm{virt}}=[\overline{M}_{0,n}(X,\beta)]^{\mathrm{virt}}.$$
\end{proposition}

\begin{corollary}[{\cite[Corollary 5.15]{manolache_VirtualPullbacks2012}}]\label{manolacheVirtualPullbacks2012-cor-5.15}
 Let $X$ and $Y$ as above, and let $\gamma\in  A^*(X)^{\otimes n}$ be any $n$-tuple of classes such that $\sum codim(\gamma_i)=vdim\overline{M}_{0,n}(X,\beta)$. Then, $I_{0,n,\tilde{\beta}}^{\tilde{X}}(p_X^*\gamma)={I}_{0,n,\beta}^{X}(\gamma)$.
\end{corollary}

\begin{remark}[{\cite[Remark 5.16]{manolache_VirtualPullbacks2012}}]\label{manolacheVirtualPullbacks2012-rem-5.16}
 The equality in the statement of proposition \ref{manolacheVirtualPullbacks2012-prop-5.14} was obtained in \cite{h} in a more general context, namely under the assumption $N_{Y/X}$ is convex with an extra minor assumption. Lai analyzes the map  $\bar{p}_X:\overline{M}_{0,n}(\tilde{X},\tilde{\beta})\rightarrow\overline{M}_{0,n}(X,\beta)$ and he uses \textit{absolute} obstruction theories (see [ibid.], Section 2). These induce a perfect relative obstruction to $\bar{p}$. Lai analyzes the normal cones of $\overline{M}_{0,n}(\tilde{X},\tilde{\beta})$ and $\overline{M}_{0,n}(X,\beta)$ and he uses the relation between them in order to obtain that $(\bar{p}_X)_*[\overline{M}_{0,n}(\tilde{X}, \tilde{\beta})]^{\mathrm{virt}}=[\overline{M}_{0,n}(X,\beta)]^{\mathrm{virt}}$ (see \cite{h} Theorem 4.11). In our language Lai's assumptions imply that $\bar{p}_X$ admits a virtual pull-back. We should stress however, that we cannot use the usual relative obstruction theories to $\mathfrak{M}_{0,n}$ in order to deduce $\bar{p}^![\overline{M}_{0,n}(X,\beta)]^{\mathrm{virt}}=[\overline{M}_{0,n}(\tilde{X},\tilde{\beta})]^{\mathrm{virt}}$. More precisely, the following diagram
\begin{equation*}  \xymatrix{\overline{M}_{0,n}(\tilde{X},\tilde{\beta})\ar[rr]^{\bar{p}}\ar[rd]_{\epsilon_{\tilde{X}}}&&\overline{M}_{0,n}(X,\beta)\ar[ld]^{\epsilon_{X}}\\
&\mathfrak{M}_{0,n}}
\end{equation*}
is not commutative. To see this, one can take a map $(C,x_1,...,x_n, f)\in \overline{M}_{0,n}(\tilde{X},\tilde{\beta})$, with $C$ a reducible curve with two components $C_1$ and $C_2$ intersecting in one point and such that $C_1$ has no marked points and it is contracted by $p\circ f$. We have that $$\epsilon_{\tilde{X}}(C,x_1,...,x_n, f)=C$$ while $$\epsilon_X\circ\bar{p}(C,x_1,...,x_n, f)=C_2.$$ This shows that Corollary \ref{manolacheVirtualPullbacks2012-cor-4.9} does not apply to the above diagram.
\end{remark}

\begin{remark}[{\cite[Remark 5.17]{manolache_VirtualPullbacks2012}}]\label{manolacheVirtualPullbacks2012-rem-5.17}
 If $X$ is the zero-locus of a section $s\in H^0(\mathbb{P}, V)$, for some convex vector bundle $V$ on $\mathbb{P}$ and $Y$ respects the hypothesis of Proposition \ref{manolacheVirtualPullbacks2012-cor-4.9}, then the equality $(\bar{p}_X)_*[\overline{M}_{0,n}(\tilde{X}, \tilde{\beta})]^{\mathrm{virt}}=[\overline{M}_{0,n}(X,\beta)]^{\mathrm{virt}}$ follows from the ``Conjecture'' proved in \cite{K:01} as described below. In notations of [ibid.] we have
\begin{equation*}
 \bar{i}_*[\overline{M}_{0,n}(X,\beta)]^{\mathrm{virt}}=c_{top}(\mathcal{R}^0(\pi_{\mathbb{P}})_*ev_{\mathbb{P}}^*V )\cdot[\overline{M}_{0,n}(\mathbb{P},i_*\beta)]^{\mathrm{virt}}.
\end{equation*}
Again, using the isomorphism in Proposition \ref{manolacheVirtualPullbacks2012-prop-5.4} (ii) we get the same relation with blow-ups, namely,
\begin{equation*}
 \bar{j}_*[\overline{M}_{0,n}(\tilde{X},\tilde{\beta})]^{\mathrm{virt}}=c_{top}(\bar{p}^*\mathcal{R}^0(\pi_{\tilde{\mathbb{P}}})_*ev_{\tilde{\mathbb{P}}}^*V)\cdot [\overline{M}_{0,n}(\tilde{\mathbb{P}},\widetilde{i_*\beta})]^{\mathrm{virt}}.
\end{equation*}
Now, the equality follows from the projection formula.
\end{remark}

\subsection{Projective bundles}\label{manolacheVirtualPullbacks2012-sec-5.4}

\begin{definition}[{\cite[Definition 5.18]{manolache_VirtualPullbacks2012}}]\label{manolacheVirtualPullbacks2012-def-5.18}
Let $\beta\in A_1(X)$, then we call $\tilde{\beta}:=s(\beta)$ the lifting of $\beta$.
\end{definition}

\begin{remark}[{\cite[Remark 5.19, p.~42]{manolache_VirtualPullbacks2012}}]\label{manolacheVirtualPullbacks2012-rem-5.19}
\label{manolacheVirtualPullbacks2012-smirr}
Let $X$ be a convex variety, $\beta\in A_1(X)$, $W$ a rank-$(r-1)$ vector bundle on $X$, $q\in\mathbb{Z}$ and $\alpha\in A^*(\overline{M}_{0,n}(\mathbb{P}_X(V),\tilde{\beta}+qf_X))$ satisfying condition (37). Let $Z_1,...Z_k$ be the connected components of $\overline{M}_{0,n}(X,\beta)$. As $X$ is convex, by dimensional reasons $$(\bar{p}_X)_* \left(\alpha\cdot[\overline{M}_{0,n}(\mathbb{P}_X(V),\tilde{\beta}+qf_X)]^{\mathrm{virt}}\right)=\sum_{i=1}^k N_i[Z_i],$$ for some $N_i\in\mathbb{Q}$, possibly zero.
\\In particular, if $X=\mathbb{P}^1$, then $\overline{M}_{0,n}(X,\beta)$ is smooth, irreducible (and unobstructed) and therefore $$(\bar{p}_X)_*\left(\alpha\cdot[\overline{M}_{0,n}(\mathbb{P}_X(V),\tilde{\beta}+qf_X)]^{\mathrm{virt}}\right)=N[\overline{M}_{0,n}(X,\beta)]^{\mathrm{virt}}$$ for some $N\in\mathbb{Q}$.
\end{remark}

\begin{definition}[{\cite[Definition 5.20, p.~43]{manolache_VirtualPullbacks2012}}]\label{manolacheVirtualPullbacks2012-def-5.20}
\label{manolacheVirtualPullbacks2012-locconst}(i) In notations as above, we consider the locally constant function
\begin{equation*}
N_{X,W,\beta,q}(\alpha):\overline{M}_{0,n}(X,\beta)\to\mathbb{Q}
\end{equation*}
defined  by $N_{X,W,\beta,q}(\alpha)=N_i$ on $Z_i$.
\\(ii) Let $X=\mathbb{P}^1$, $d=(d_1,...,d_{r-1})\in\mathbb{Z}^{r-1}$, $V=\mathcal{O}\oplus \mathcal{O}_{\mathbb{P}^1}(d_1)\oplus...\oplus\mathcal{O}_{\mathbb{P}^1}(d_{r-1})$,  and $k=(k_1,...,k_n)\in \mathbb{N}^n$. Let $\xi_{\mathbb{P}_X(V)}=c_1(\mathcal{O}_{\mathbb{P}_X(V)}(1))$ and assume that $\alpha=ev^*_i\xi^{k_i}$ satisfies the dimension condition (37). We define $N(q,d,k)\in \mathbb{Q}$ by the formula $$(\bar{p}_X)_*\left(\alpha\cdot[\overline{M}_{0,n}(\mathbb{P}_X(V),\tilde{1}+qf_X)]^{\mathrm{virt}}\right)=N(q,d, k)[\overline{M}_{0,n}(X,1)].$$
\end{definition}

\begin{remark}[{\cite[Remark 5.21]{manolache_VirtualPullbacks2012}}]\label{manolacheVirtualPullbacks2012-rem-5.21}
Let $X$ and $Y$ be smooth projective varieties, let $f:Y\to X$ be a morphism and $\beta_Y\in A_1(Y)$ such that $f_*\beta_Y=\beta$. Let $W$ be a vector bundle on $X$. Then there exists an induced map $h:\mathbb{P}_Y(f^*V)\to\mathbb{P}_X(V)$. This induces a map $\bar{h}:\overline{M}_{0,n}(\mathbb{P}_Y(f^*V),\tilde{\beta}_Y+qf_Y)\to \overline{M}_{0,n}(\mathbb{P}_X(V),\tilde{\beta}+qf_X)$
\end{remark}

\begin{proposition}[{\cite[Proposition 5.22, p.~43]{manolache_VirtualPullbacks2012}}]\label{manolacheVirtualPullbacks2012-prop-5.22}
\uses{manolacheVirtualPullbacks2012-prop-5.4, manolacheVirtualPullbacks2012-def-5.20, manolacheVirtualPullbacks2012-thm-4.1, manolacheVirtualPullbacks2012-cor-4.9}
\label{manolacheVirtualPullbacks2012-projb} Let $X$ be smooth convex projective varieties and let $f:Y\to X$ be a morphism of smooth projective varieties such that $$f_*:A_1(Y)\to A_1(X)$$ is injective. Let $W$ be a vector bundle on $X$ and $$\alpha\in A^*(\overline{M}_{0,n}(\mathbb{P}_X(V),\tilde{\beta}+qf_X))$$ which satisfies the dimension condition (37). Then
\\(i) $(\beta_Y, f^*W,q, \bar{h}^*\alpha)$ satisfies condition 37.
\\(ii)Let $Z_1,...,Z_k$ be the connected components of $\overline{M}_{0,n}(X,\beta)$. For any $i\in\{1,...,k\}$ let $N_{X,W,\beta,q}(\alpha)|_{Z_i}$ be the number from definition \ref{manolacheVirtualPullbacks2012-def-5.20} which corresponds to $Z_i$. Then, we have an equality
\begin{equation*}

(\bar{p}_Y)_*\left((\bar{h}^*\alpha)\cdot[\overline{M}_{0,n}(\mathbb{P}_{Y}(f^*V),\tilde{\beta_Y}+qf_Y)]^{\mathrm{virt}}\right)=\sum_{i=1}^kN_{X,W,\beta,q}(\alpha)|_{Z_i}[Z_i]
\end{equation*}
\end{proposition}

\begin{corollary}[{\cite[Corollary 5.23, p.~44]{manolache_VirtualPullbacks2012}}]\label{manolacheVirtualPullbacks2012-cor-5.23}
\uses{manolacheVirtualPullbacks2012-prop-5.22}
\label{manolacheVirtualPullbacks2012-convexbdl} We follow the notations of Proposition \ref{manolacheVirtualPullbacks2012-prop-5.22}. Let $X$ be a homogeneous space, $Y$ a rational curve. Let $$\alpha:=\prod_{i=1}^nev_i^*\xi^{k_i}\in A^*(\overline{M}_{0,n}(\mathbb{P}_{X}(V),\tilde{\beta}+qf_{X}))$$ satisfy condition (37). Suppose $f^*V\simeq \mathcal{O}\oplus\mathcal{O}(d_1)\oplus...\oplus\mathcal{O}(d_{r-1})$. Then $$(\bar{p}_{X})_* \left(\alpha\cdot[\overline{M}_{0,n}(\mathbb{P}_{X}(V),\tilde{\beta}+qf_{X})]^{\mathrm{virt}}\right)=N(q,d, k)[\overline{M}_{0,n}(Y,\beta)]^{\mathrm{virt}}.$$
\end{corollary}

\begin{corollary}[{\cite[Corollary 5.24, p.~45]{manolache_VirtualPullbacks2012}}]\label{manolacheVirtualPullbacks2012-cor-5.24}
\uses{manolacheVirtualPullbacks2012-prop-5.22, manolacheVirtualPullbacks2012-cor-5.23}
\label{manolacheVirtualPullbacks2012-jump} Let $d=(d_1,...,d_{r-1})$, $e=(e_1,...,e_{r-1})$ with $d_i,\ e_i\geq 0$, $\forall\ 0\leq i\leq r$ and $\sum_{i=1}^{r-1} d_i=\sum_{i=1}^{r-1} e_i$. Then $N(q,d,k)=N(q,e,k)$.
\end{corollary}

\begin{setting}[{\cite[Setting 5.25, p.~45]{manolache_VirtualPullbacks2012}}]\label{manolacheVirtualPullbacks2012-setting-5.25}
\label{manolacheVirtualPullbacks2012-notatii}We consider a homogeneous space space $\mathbb{P}$ and $g:X\to\mathbb{P}$ be a closed embedding of a smooth projective variety $X$ in $\mathbb{P}$. Let $V$ be a vector bundle on $X$ such that there exists a vector bundle $W\oplus \mathbb{C}_{\mathbb{P}}$ on $\mathbb{P}$ with $V=g^*(W\oplus \mathbb{C}_{\mathbb{P}})$ and let $p:\mathbb{P}_X(V)\to X$ be the associated projective bundle. In notations as above we have an induced map $\bar{p}_X:\overline{M}_{0,n}(\mathbb{P}_{X}(V),\tilde{\beta}+qf_X)\to\overline{M}_{0,n}(X,\beta)$.
\end{setting}

\begin{definition}[{\cite[Definition 5.26]{manolache_VirtualPullbacks2012}}]\label{manolacheVirtualPullbacks2012-def-5.26}
 Let $V$ be a vector bundle on a smooth projective variety $X$ and let $i:\mathbb{P}^1\to X$ be a closed embedding of a projective line in $X$. Let $i^*V\simeq\mathcal{O}(d_1)\oplus...\oplus\mathcal{O}(d_{r-1})$ for some $d_1,...,d_{r-1}\in\mathbb{Z}$. We say that $i^*V$ is positive if $d_i\geq0$, $\forall\ 0\leq i\leq r-1$.
\end{definition}

\begin{corollary}[{\cite[Corollary 5.27, p.~45]{manolache_VirtualPullbacks2012}}]\label{manolacheVirtualPullbacks2012-cor-5.27}
\uses{manolacheVirtualPullbacks2012-prop-5.22, manolacheVirtualPullbacks2012-cor-5.23, manolacheVirtualPullbacks2012-cor-5.24}
\label{manolacheVirtualPullbacks2012-genbdl} In notations as in \ref{manolacheVirtualPullbacks2012-setting-5.25}, let $$\alpha:=\prod_{i=1}^nev_i^*\xi^{k_i}\in A^*(\overline{M}_{0,n}(\mathbb{P}_X(V),\tilde{\beta}+qf_X))$$ satisfy condition (37). Suppose moreover that the restriction of $V$ to any curve of class $\beta$ is positive. Then we have that $$(\bar{p}_X)_* \left(\alpha\cdot[\overline{M}_{0,n}(\mathbb{P}_X(V),\tilde{\beta}+qf_X)]^{\mathrm{virt}}\right)=N(q,d, k)[\overline{M}_{0,n}(X,\beta)]^{\mathrm{virt}},$$
for some $d=(d_1,...,d_n)$ which satisfies $\sum_id_i=\beta\cdot c_1(V)$.
\end{corollary}

\begin{remark}[{\cite[Remark 5.28]{manolache_VirtualPullbacks2012}}]\label{manolacheVirtualPullbacks2012-rem-5.28}
 In Corollary \ref{manolacheVirtualPullbacks2012-cor-5.27} we have shown that we can compute GW invariants of projective bundles in terms of GW invariants of the base $X$ and GW invariants of a projective bundle over $\mathbb{P}^1$. The latter can be analyzed using toric methods (see \cite{e1}, \cite{e2}). More precisely, we can compute in this way GW-invariants of $\mathbb{P}_X(V)$ with at least $\mathrm{vdim} \overline{M}_{0,n}(X,\beta)$ insertions that are pull-backs from $X$.
\end{remark}

\subsection{Costello's push-forward formula}\label{manolacheVirtualPullbacks2012-sec-5.5}

\begin{proposition}[{\cite[Proposition 5.29, p.~46]{manolache_VirtualPullbacks2012}}]\label{manolacheVirtualPullbacks2012-prop-5.29}
\uses{manolacheVirtualPullbacks2012-thm-4.1}
\label{manolacheVirtualPullbacks2012-costello} Under the assumptions above, if moreover $g$ is projective, then $f_*[F]^{\mathrm{virt}}=d[G]^{\mathrm{virt}}$.
\end{proposition}

\begin{remark}[{\cite[Remark 5.30]{manolache_VirtualPullbacks2012}}]\label{manolacheVirtualPullbacks2012-rem-5.30}
 In \cite{c}, Theorem 5.0.1 $g$ is not assumed to be projective. We impose this condition in order to be able to push-forward along $g$ and apply Theorem \ref{manolacheVirtualPullbacks2012-thm-4.1}. If $\mathfrak{M}_1$ and $\mathfrak{M}_2$ are DM-stacks then it is enough to assume $g$ is proper. The proof of Proposition \ref{manolacheVirtualPullbacks2012-prop-5.29} applies unchanged to this case.
\\More generally, Theorem 5.0.1 in \cite{c} follows from the proof of Theorem \ref{manolacheVirtualPullbacks2012-thm-4.1} (i) which is very similar to Costello's proof.
\end{remark}
